\section{The ordinary resolution and its filtered pairing}
\label{app:resolution}

We use a symmetric monoidal stable model of spectra with the usual
pushout-product axiom and tensors over based spaces. Such a model is
provided by topological symmetric spectra:
\citet[Section~6.2, Theorems~6.3.8 and~6.4.1, and
Corollary~6.4.2]{HoveyShipleySmith1998v2}; its comparison with ordinary
spectra follows from Theorems~6.3.8 and~4.3.2 there, in the cited
preprint version.
All smash products
are derived; when point-set maps are used, their inputs are cofibrant.
Rectify the coherent associative structure on $BP$ as an $S$-algebra by
\citet[Part~II, Corollary~3.6 and the associative convention comparison
in Section~4]{ElmendorfKrizMandellMay1997}. The equivalence for associative
algebras in \citet[Main Theorem and Theorem~5.1(ii)]{Schwede2001} then
places it in this same category of topological symmetric spectra.
Choose a cofibrant associative symmetric ring representative using
\citet[Theorem~12.1(iv)--(v)]{MandellMaySchwedeShipley2001}; its
underlying spectrum is then stably cofibrant. Smash with such a spectrum
preserves stable equivalences by Proposition~12.3 there, so the following
point-set smash powers compute derived smash powers.
The comparison uses the positive stable model structure, whose weak
equivalences agree with the ordinary stable ones; we use the ordinary
stable structure here. With this strictly unital associative model, write
\[
 T^n(Z)=BP^{\wedge(n+1)}\wedge Z.
\]
The cofaces insert units and the codegeneracies multiply adjacent $BP$
factors.  Associativity and the unit identities give a cosimplicial
spectrum.  Write $K_n(Z)=\operatorname{Tot}_n T^\bullet(Z)$ for its
derived ordinary partial totalization, set $K_{-1}(Z)=*$, and put
$K(Z)=\operatorname{Tot}T^\bullet(Z)$.  Ordinary totalization here uses
all ordinal maps, including codegeneracies.

\subsection{Convergence and the edge}\label{app:convergence}

\begin{proof}[Proof of Proposition~\ref{fw:anss}]
We first identify the ordinary totalization tower exactly.  Let
$\mathcal P_0([n])$ be the poset of nonempty subsets of
$\{0,\ldots,n\}$.  The functor
\[
 u_n:\mathcal P_0([n])\longrightarrow\Delta_{\leq n},
 \qquad S\longmapsto[|S|-1],
\]
sends an inclusion to the corresponding order-preserving injection.
It is homotopy-initial.  To prove this, for $m\leq n$ identify
$(u_n\downarrow[m])$ with the nonempty face poset of the simplicial
complex $L_{n,m}$ whose vertices are $(i,j)$, $0\leq i\leq n$,
$0\leq j\leq m$, and whose simplices have strictly increasing first
coordinates and weakly increasing second coordinates.  Its nerve is
the barycentric subdivision of $L_{n,m}$.

The complex $L_{n,m}$ is contractible for $n\geq m$, by induction on
$n$.  The case $n=m=0$ is a point.  For the inductive step start with
the cone with vertex $(n,m)$ on the entire subcomplex $L_{n-1,m}$.
Adjoin the cones with vertices $(n,j)$, $0\leq j<m$, on their bases
$L_{n-1,j}$.  Two distinct vertices with first coordinate $n$ cannot
belong to a simplex.  Each new cone therefore intersects the preceding
union in exactly its base.  That base is contractible by induction,
since $n-1\geq j$.  Attachment along this simplicial subcomplex is a
homotopy pushout and preserves contractibility.  This proves the claim.

Homotopy-initiality, applied to the cosimplicial unit resolution, now
gives, naturally in $n$ and $Z$,
\begin{equation}
 K_n(Z)\simeq
 \operatorname*{holim}_{\varnothing\ne S\subseteq\{0,\ldots,n\}}
       BP^{\wedge|S|}\wedge Z.
 \label{fw:finite-cube-comparison}
\end{equation}
This is the standard cubical description of the partial Amitsur
totalizations; compare \citet[Propositions~2.11 and~2.14]{MathewNaumannNoel2017}.
Indeed derived partial totalization is the homotopy limit over
$\Delta_{\leq n}$: the restricted standard simplex is a Reedy-cofibrant
resolution of the constant point diagram.  Adjoining the empty vertex
$Z$ on the right of Equation~\eqref{fw:finite-cube-comparison} gives
the $(n+1)$-fold cube formed by smashing copies of the unit
$S_{(p)}\to BP$.  If $\mathcal F$ is its fiber, iterated fibers and
exactness of smash identify its total fiber with
$\mathcal F^{\wedge(n+1)}\wedge Z$.  Thus
\begin{equation}
 \operatorname{fib}(Z\longrightarrow K_n(Z))
       \simeq\mathcal F^{\wedge(n+1)}\wedge Z.
 \label{fw:tot-fiber-index}
\end{equation}
The comparisons commute with restriction from $n+1$ to $n$: on the
fibers this is the map induced by the last factor
$\mathcal F\to S_{(p)}$.  There is no suspension in
Equation~\eqref{fw:tot-fiber-index}.

The unit induces an isomorphism on $\pi_0$, and $\pi_1BP=0$, so
$\mathcal F$ has no homotopy in degrees at most zero.  If $Z$ is
$c$-connective, $\mathcal F^{\wedge s}\wedge Z$ is at least
$(c+s)$-connective.  Equation~\eqref{fw:tot-fiber-index} therefore
identifies the coaugmentation $Z\to K(Z)$ as an equivalence, and
identifies the relative totalization tower as
\begin{equation}
 G^s(Z):=\operatorname{fib}(K(Z)\longrightarrow K_{s-1}(Z))
        \simeq\mathcal F^{\wedge s}\wedge Z
        \quad(s\geq0).
 \label{fw:relative-tower}
\end{equation}
Its homotopy groups vanish in each fixed stem for sufficiently large
$s$.  The exact couple consequently has the asserted abutment, no
inverse-limit ambiguity, and finite filtration in each stem.

For completeness, the successive layer is
\begin{equation}
 \operatorname{cofib}(G^{s+1}(Z)\longrightarrow G^s(Z))
 \simeq BP\wedge\mathcal F^{\wedge s}\wedge Z
 \simeq\Sigma^{-s}BP\wedge\overline{BP}^{\wedge s}\wedge Z,
 \label{fw:tot-layer-index}
\end{equation}
where $\overline{BP}=\operatorname{cofib}(S_{(p)}\to BP)$.
This is also $\operatorname{fib}(K_s(Z)\to K_{s-1}(Z))$, with the
usual normalization and the displayed $s$ desuspensions.

The iterated extended terms compute the cobar complex.  Indeed
$BP\wedge BP$, as a $BP$-module, is a wedge of suspended free
$BP$-modules with basis the polynomial monomials in the $t_i$.
Iteration identifies their homotopy with the iterated tensors of
$\Gamma$ over $R$.  There are no derived tensor corrections, whether
or not $BP_*Z$ is flat over $R$.  Cofaces give the cobar differential;
the codegeneracy splittings give its normalized form.  The exact-couple
grading gives the displayed differential shift.  Finally the first
map is $Z\to BP\wedge Z$.  Its image is exactly the permanent
filtration-zero part by the finite-filtration convergence just proved.
\end{proof}

\subsection{The filtered binary pairing}\label{app:pairing}

\begin{proof}[Proof of Proposition~\ref{fw:anss-pairing}]
We construct maps
\[
 G^s(Z_1)\wedge G^t(Z_2)\longrightarrow G^{s+t}(Z_3)
\]
which induce the stated cobar pairing and recover the actual map $f$.

\smallskip\noindent\emph{The box pairing.}
We use Batanin's cosimplicial box product in the form of
\citet[Definition~2.1, Proposition~2.3, and Remark~2.4]{McClureSmith2004};
they attribute its Kan-extension description to Cordier and Porter.
We construct the pairing before any fibrant replacement and then transport
the entire pairing. For cosimplicial spectra define $A\mathbin\square B$
as the left Kan extension of $(a,b)\mapsto A^a\wedge B^b$ along ordinal
concatenation
\[
 \Delta\times\Delta\longrightarrow\Delta,
 \qquad([a],[b])\longmapsto[a+b+1].
\]
Equivalently, in degree $n$ it is the zigzag pushout of the terms
$A^a\wedge B^b$ with $a+b=n$, with neighboring terms identified by
the maps
\[
 A^a\wedge B^b\xrightarrow{d^{a+1}\wedge1}
 A^{a+1}\wedge B^b,
 \qquad
 A^a\wedge B^b\xrightarrow{1\wedge d^0}
 A^a\wedge B^{b+1},
 \qquad a+b=n-1.
\]
The same definition applies to cosimplicial based spaces.
For the original unit resolutions, concatenating the two $BP$ blocks
and applying $f$ once gives maps
\[
 T^a(Z_1)\wedge T^b(Z_2)\longrightarrow T^{a+b+1}(Z_3).
\]
They are natural in both ordinals: cofaces and codegeneracies within a
block insert a unit or multiply adjacent factors within that block.
Kan-extension adjunction therefore gives an actual cosimplicial map
\begin{equation}
 m:T(Z_1)\mathbin\square T(Z_2)\longrightarrow T(Z_3).
 \label{fw:original-box-pairing}
\end{equation}
On the degree-$n$ summand $a+b=n$ in the pushout description, this map
multiplies the last factor of the first $BP$ block with the first factor
of the second block.  This is the single overlapping slot.  The two
descriptions agree by applying the codegeneracy that identifies the
two junction vertices of $[a+b+1]$.

\smallskip\noindent\emph{Deriving the pairing.}
We record the model-category check that permits us to derive this
construction.  The box product is a Quillen bifunctor for the Reedy
model structures.  Let $h_a([n])=\Delta([a],[n])_+$, and let
$\partial h_a$ be the subdiagram of noninjective maps.  These are the
representables and boundaries defining the generating Reedy cells.
Indeed, the boundary of the covariant representable consists of maps
factoring through a nonidentity inverse map out of $[a]$, hence of
noninjective maps by the unique surjection--injection factorization in
$\Delta$. This representable boundary is distinct from the latching
subobject in target degree $n$, which consists of maps factoring through
a proper injection into $[n]$.
Kan extension gives $h_a\square h_b=h_{a+b+1}$.  In the relative box
boundary
\[
 \frac{h_{a+b+1}}
 {\operatorname{im}(\partial h_a\square h_b)
       \cup\operatorname{im}(h_a\square\partial h_b)},
\]
the remaining maps are exactly the monotone maps injective on each of
the two blocks.  Such a map is either globally injective or identifies
only the two junction vertices. A map that is not surjective onto its
target lies in that target degree's latching subobject. The only
remaining surjections are the junction identification onto $[a+b]$
and the identity of $[a+b+1]$. Thus the junction-identifying and globally
injective non-latching cells occur in degrees $a+b$ and $a+b+1$,
respectively.  Attach the junction cell first;
every further collapse of it is already in the boundary, since it
collapses two vertices within one block.  Attach the globally injective
cell second: its junction codegeneracy lands in the first cell, and
its other codegeneracies land in the boundary. This gives a two-cell
Reedy filtration of the relative box boundary. For two generating
Reedy cofibrations, each of these attachments is tensored with the
ordinary spectral pushout-product of their underlying generating
cofibrations. The spectral pushout-product axiom therefore makes both
attachments cofibrations, and makes them trivial cofibrations if either
input is trivial. This proves the Reedy pushout-product axiom, including
the acyclic case, for generating cofibrations.  Closure under pushouts,
composition, and retracts proves the assertion for all cofibrations.

Choose Reedy-cofibrant replacements $A_i\to T(Z_i)$, $i=1,2$, that
are trivial fibrations, and trivial cofibrations $A_i\to R_i$ with
$R_i$ Reedy fibrant.  For the target choose a Reedy fibrant replacement
$T(Z_3)\to R_3$ directly.  Equation~\eqref{fw:original-box-pairing}
gives $A_1\square A_2\to R_3$.  By the preceding Quillen check,
$A_1\square A_2\to R_1\square R_2$ is a trivial Reedy cofibration.
The lifting axiom extends the map to
\begin{equation}
 m_R:R_1\square R_2\longrightarrow R_3.
 \label{fw:derived-box-pairing}
\end{equation}
The space of such extensions is contractible up to the usual relative
homotopy.  This transports the original pairing; it places no condition
on the literal $BP$ slots of the fibrant replacements.
All resulting tower and spectral-sequence pairings are transported along
the equivalences with the original diagrams.  In particular, when some
$Z_i$ coincide, their source and target replacements need not be identical.

\smallskip\noindent\emph{The filtered diagonal.}
We use the simplex subdivision map given in
\citet[Lemma~3.6]{McClureSmith2004}, credited to Grayson in their
Remark~3.7. We record its formula to verify the relative filtration. For
$t=(t_0,\ldots,t_n)\in\Delta^n$, put
\[
 I_i=\left[\sum_{j<i}t_j,\sum_{j\leq i}t_j\right],\qquad
 x_i=2\,\operatorname{length}(I_i\cap[0,1/2]),\qquad
 y_i=2t_i-x_i.
\]
For a cut index $k$, $x$ is supported in $0,\ldots,k$ and $y$ in
$k,\ldots,n$.  The cut region is a copy of
$\Delta^k\times\Delta^{n-k}$, with inverse $t=(x+y)/2$.
Its boundary seams identify the last face of the first factor with
the first face of the second factor, precisely as in the box pushout.
Thus the cut defines
\[
 \delta:\Delta^\bullet_+\longrightarrow
              \Delta^\bullet_+\square\Delta^\bullet_+.
\]
It is natural under all monotone maps: they aggregate consecutive
intervals, and aggregation commutes with the cut.  In particular this
is a map of ordinary cosimplicial objects, including codegeneracies.

Define the relative simplex object
\[
 P_s^\bullet=\Delta^\bullet_+/\operatorname{sk}_{s-1}
                      \Delta^\bullet_+\qquad(s\geq0),
\]
where the skeleton for $s=0$ is the basepoint.  The simplex, its
skeleta, and these quotients are Reedy cofibrant: their latching maps
are the corresponding boundary inclusions or identities.  Consequently
\begin{equation}
 G^s(Z_i)=\operatorname{Map}_\Delta(P_s,R_i)
 \label{fw:relative-end-model}
\end{equation}
is a derived relative end and the actual fiber of the fibration
$\operatorname{Tot}R_i\to\operatorname{Tot}_{s-1}R_i$.
The degree-$n$ box summands have dimensions $k$ and $n-k$.
If $n<s+t$, at least one is below its respective skeleton cutoff.
Naturality applies the same observation to every face of a larger
simplex, including faces with missing ambient vertices.  Hence the
cut descends to maps
\begin{equation}
 \delta_{s,t}:P_{s+t}\longrightarrow P_s\square P_t.
 \label{fw:relative-cut}
\end{equation}
Enriched evaluation, followed by $m_R$ and precomposition with
$\delta_{s,t}$, gives the actual relative pairing
\begin{align*}
 \operatorname{Map}_\Delta(P_s,R_1)\wedge
 \operatorname{Map}_\Delta(P_t,R_2)
 &\longrightarrow
 \operatorname{Map}_\Delta(P_s\square P_t,R_1\square R_2)\\
 &\longrightarrow
 \operatorname{Map}_\Delta(P_{s+t},R_3).
\end{align*}
It commutes with the tower inclusions in both inputs.  By
Equation~\eqref{fw:relative-tower}, these are the exact relative fibers
of the unit-resolution filtration, with filtration indices $s,t,s+t$.

\smallskip\noindent\emph{The cobar and homotopy pairings.}
Put $L_s=\operatorname{cofib}(G^{s+1}\to G^s)$.  The relative pairing
induces $L_s(Z_1)\wedge L_t(Z_2)\to L_{s+t}(Z_3)$: raising either
input filtration by one raises the output filtration by one.
The boundary formula for the product of relative cells is
$\partial(a\cdot b)=\partial a\cdot b+(-1)^{|a|}a\cdot\partial b$.
It gives a pairing of the exact couples and the Leibniz rule on every
derived page.  On the layer of dimension $s+t$, the only contributing
cut cell has front and back dimensions $s,t$.  The inverse
$t=(x+y)/2$ gives its usual oriented front-face/back-face cup map,
with coefficient one.  The normalized cobar identification therefore
gives exactly the asserted $E_2$ cup pairing.

We finally verify the homotopy target without assuming that a replacement
preserves units.  For cofibrant $Z_i$, the diagram
$\Delta^\bullet_+\wedge Z_i$ is Reedy cofibrant and maps levelwise
equivalently to the constant diagram $Z_i$.  Its original all-unit map
to $T(Z_i)$ lifts to a map
$\alpha_i:\Delta^\bullet_+\wedge Z_i\to A_i$ along the trivial
fibration $A_i\to T(Z_i)$.  The composite of $\alpha_i$ with
$A_i\to R_i$ represents the coaugmentation.  The box product of the
$\alpha_i$, precomposed with $\delta$, lands in $A_1\square A_2$.
Since $m_R$ extends the
original map on that object, the composite is exactly the target
coaugmentation composed with $f$, transported through $T(Z_3)\to R_3$.
On totalizations the pairing therefore satisfies
\[
 m_{\mathrm{Tot}}\circ(e_1\wedge e_2)=e_3\circ f.
\]
The coaugmentations $e_i:Z_i\to K(Z_i)$ are the equivalences proved
above.  Thus the limit pairing is the actual homotopy map $f$, and
the $E_\infty$ pairing is its associated-graded pairing for the
filtration by the images of $\pi_*G^s\to\pi_*K$.

For the final assertion, the coefficient pairing is $R$-bilinear and
sends the two units to the unit.  Every element of $R/I$ is a scalar
multiple of that unit.  The pairing is consequently
$(r\bmod I,r'\bmod I)\mapsto rr'\bmod I$.  This uses neither flatness
of the quotient nor associativity of the target map.
\end{proof}
